% =====================================================================
%  THÈME 1 — Pondération des équations chimiques
%  Niveau FACILE — Corrigé
% =====================================================================
\documentclass[11pt,a4paper]{article}
\usepackage[utf8]{inputenc}
\usepackage[T1]{fontenc}
\usepackage{lmodern}
\usepackage[french]{babel}
\usepackage{amsmath,amssymb,mathtools}
\usepackage{icomma}
\usepackage[version=4]{mhchem}
\usepackage{siunitx}
\sisetup{output-decimal-marker={,},per-mode=symbol}
\usepackage{xcolor}
\usepackage{tikz}
\usepackage[most]{tcolorbox}
\usepackage{enumitem}
\usepackage{array}
\usepackage{booktabs}
\usepackage{fancyhdr}
\usepackage[a4paper,margin=2.1cm,headheight=26pt,headsep=8pt]{geometry}

\definecolor{primary}{RGB}{142,93,168}
\definecolor{primarydark}{RGB}{82,45,108}
\definecolor{corrcol}{RGB}{176,110,20}

\pagestyle{fancy}\fancyhf{}
\fancyhead[L]{\small\textcolor{primarydark}{\textbf{Fiche 5} — Thème 1 : Pondération}}
\fancyhead[R]{\small\textcolor{corrcol}{\textsc{Corrigé — Facile}}}
\fancyfoot[C]{\small\thepage}
\renewcommand{\headrulewidth}{0.8pt}
\renewcommand{\headrule}{\hbox to\headwidth{\color{primary}\leaders\hrule height \headrulewidth\hfill}}

\newtcolorbox[auto counter]{cor}[1][]{enhanced,breakable,boxrule=1pt,arc=3pt,
  left=3mm,right=3mm,top=2.4mm,bottom=2.4mm,colback=corrcol!6,colframe=corrcol,
  fonttitle=\bfseries,coltitle=white,
  title={Corrigé~\thetcbcounter\ \ifx\relax#1\relax\relax\else\textmd{--- \itshape #1}\fi}}

\newcommand{\rep}[1]{\par\smallskip\noindent\textbf{\color{corrcol}$\Rightarrow$ #1}}

\setlength{\parindent}{0pt}
\setlength{\parskip}{3pt}

\begin{document}

\begin{center}
{\LARGE\bfseries\color{primarydark}Thème 1 — Pondération}\\[1pt]
{\large\color{corrcol}\textbf{Corrigé — Niveau Facile}}\\
{\color{primary}\rule{0.6\linewidth}{1pt}}
\end{center}

\begin{cor}[synthèse de l'ammoniac]
\textbf{a)} On équilibre l'azote puis l'hydrogène. À droite \ce{NH3} : pour avoir 2 N, il faut
$2\,\ce{NH3}$, ce qui donne $2\times3=6$ atomes H ; à gauche il faut donc $3\,\ce{H2}$ ($=6$ H) :
\[ \boxed{\ce{N2 + 3 H2 -> 2 NH3}} \]
\emph{Vérif :} N $2=2$ ; H $6=6$. \checkmark
\rep{b) Somme $=1+3+2=\boxed{6}$.}
\end{cor}

\begin{cor}[combustion du méthane]
\textbf{a)} Carbone : 1 C $\to 1\,\ce{CO2}$. Hydrogène : 4 H $\to 2\,\ce{H2O}$ ($2\times2=4$).
Oxygène à droite : $1\times2+2\times1=4$ atomes O $\Rightarrow 2\,\ce{O2}$ à gauche :
\[ \boxed{\ce{CH4 + 2 O2 -> CO2 + 2 H2O}} \]
\rep{b) De chaque côté : \textbf{4 atomes d'oxygène} ($2\,\ce{O2}=4$ à gauche ; $2+2=4$ à droite). \checkmark}
\end{cor}

\begin{cor}[oxydation du fer]
Fer : \ce{Fe2O3} impose 2 Fe $\Rightarrow 2\,\ce{Fe}$. Oxygène : 3 O à droite face à des paquets de 2 ;
ppcm$(2,3)=6$, d'où $3\,\ce{O2}$ (6 O) à gauche et $2\,\ce{Fe2O3}$ (6 O) à droite. On re-corrige le fer :
$2\,\ce{Fe2O3}=4$ Fe $\Rightarrow 4\,\ce{Fe}$ :
\[ \boxed{\ce{4 Fe + 3 O2 -> 2 Fe2O3}} \]
\emph{Vérif :} Fe $4=4$ ; O $6=6$. \checkmark
\rep{Somme $=4+3+2=\boxed{9}$.}
\end{cor}

\begin{cor}[repérer une erreur de pondération]
\textbf{a) Inventaire de « \ce{2 Al + O2 -> Al2O3} ».}
\begin{center}\small
\begin{tabular}{lcc}
\toprule
Élément & Réactifs & Produits\\ \midrule
\ce{Al} & $2$ & $2$ \quad\checkmark\\
\ce{O}  & $2$ (\ce{O2}) & $3$ (\ce{Al2O3}) \quad\textcolor{red}{$\neq$}\\
\bottomrule
\end{tabular}
\end{center}
L'oxygène n'est \textbf{pas} équilibré ($2\neq3$) : l'équation proposée est \textbf{fausse}.

\textbf{b)} On équilibre l'oxygène par ppcm$(2,3)=6$ : $3\,\ce{O2}$ (6 O) à gauche, $2\,\ce{Al2O3}$
(6 O) à droite ; puis l'aluminium : $2\,\ce{Al2O3}=4$ Al $\Rightarrow 4\,\ce{Al}$ :
\[ \boxed{\ce{4 Al + 3 O2 -> 2 Al2O3}} \]
\rep{Somme $=4+3+2=\boxed{9}$.}
\end{cor}

\begin{cor}[nommer le type de réaction]
\begin{enumerate}[label=\arabic*., leftmargin=1.8em,topsep=1pt,itemsep=1pt]
  \item \ce{AgNO3 + KBr -> AgBr(s) + KNO3} : formation d'un solide peu soluble (\ce{AgBr})
        $\Rightarrow$ \textbf{précipitation}.
  \item \ce{CaCO3 -> CaO + CO2} : une seule espèce se scinde en deux $\Rightarrow$ \textbf{décomposition}.
  \item \ce{C2H4 + H2 -> C2H6} : addition de \ce{H2} sur un alcène $\Rightarrow$ \textbf{hydrogénation}.
  \item \ce{H2SO4 + 2 NaOH -> Na2SO4 + 2 H2O} : transfert de protons entre un acide et une base
        $\Rightarrow$ \textbf{acide--base}.
  \item \ce{4 Fe + 3 O2 -> 2 Fe2O3} : le fer perd des électrons, l'oxygène en gagne
        $\Rightarrow$ \textbf{redox} (oxydoréduction).
\end{enumerate}
\rep{$1=$ précipitation ; $2=$ décomposition ; $3=$ hydrogénation ; $4=$ acide--base ; $5=$ redox.}
\end{cor}

\end{document}
